\documentclass[11pt,a4paper]{article}
\usepackage[T1]{fontenc}
\usepackage[utf8]{inputenc}
\usepackage[english]{babel}
\usepackage{amsmath,amssymb}
\usepackage[margin=2cm]{geometry}
\setlength{\parindent}{0pt}
\setlength{\emergencystretch}{3em}
\allowdisplaybreaks

\title{About \texttt{false\_book\_partial\_derivatives.tex}}
\author{}
\date{}
\begin{document}
\maketitle
This note explains what the test data in \texttt{false\_book\_partial\_derivatives.tex} is used for. The file is read by \texttt{test\_multivariable\_calculus.c}.

\begin{enumerate}
\item \textbf{Purpose.} The file holds 24 false statements. The first eight are results printed wrong in the collections of solved exercises: a result of a book is data, and when it is wrong it is a false statement that must not be proved. The other sixteen are right results of \texttt{book\_partial\_derivatives.tex} changed by little.

\item \textbf{The printed errors.} Item 1: Eleuteri 3.4.2, the gradient of $\ln(1 + \sqrt{x^{2} + 2y})$ at $(2, 0)$ is $(\frac{1}{3}, \frac{1}{6})$, not $(\frac{2}{3}, \frac{1}{3})$. Item 2: UBC 2.5.2(b), the numerator is $-2x^{2} + y^{2} + z^{2}$. Item 3: Pagliarini 11, $f_{y} = e^{x + y}(x^{2} + y + 1)$. Item 4: Eleuteri 5.1.11, the Hessian matrix at $(\frac{1}{2}, \frac{1}{3})$. Item 5: Eleuteri 5.1.2, the Hessian matrix at the origin is $\begin{pmatrix} -4 & 4 \\ 4 & -4 \end{pmatrix}$. Item 6: Eleuteri 5.1.19(11), the critical point is $(0, \frac{1}{4})$, not $(0, 4)$. Item 7: Perfetti 12.7, the critical point is $(1, 1)$. Item 8: UBC 4.1.5, the last entry is $6$. Each was decided by a third computation (\texttt{photos/dati\_esercizi/analisi\_funzioni\_due\_variabili.md}, section 6), and for some on the image of the page.

\item \textbf{Expected result.} Every statement is called false, with a counterexample.

\item \textbf{How to run it.} From the top folder of the repository:
\begin{verbatim}
cc -std=c11 -O2 -Wall -Wextra -o /tmp/test_multivariable_calculus \
    tests/multivariable_calculus/test_multivariable_calculus.c && \
    /tmp/test_multivariable_calculus
\end{verbatim}
\end{enumerate}
\end{document}
